\originalchapter{22.02 作业九:贝塔谱与中子散射}
题源: Spring 2012 Problem Set 9, 2012年5月3日. 全部为编者独立解答, 核质量数值版本 AME2020.
\begin{exercise}\label{ex:2202-9-1}
原题 1. 完成太阳聚变支路: (a) $^1$H+$^1$H$\to\cdots$, (b) $^1$H+$\cdots\to{}^3$He+$\gamma$ (非 beta 过程), (c) $^1$H+$^3$He$\to{}^4$He+$\cdots+\cdots$.
\end{exercise}
\begin{solution}
依核子数、电荷和轻子数守恒: (a) d+$e^++\nu_e$; (b) 入射缺 d; (c) $e^++\nu_e$, 为 hep 支路. (b) 两体合并辐射 gamma, (a),(c) 质子转中子需弱作用. (c) 不应换成常见两 $^3$He 的 pp 主链反应.
\end{solution}
\begin{exercise}\label{ex:2202-9-2}
原题 2. 激光原子阱中静止 $^8$He 经 beta 衰变, 写反应及产物动能.
\end{exercise}
\begin{solution}
$^8$He$\to{}^8$Li+$e^-+\bar\nu_e$, 基态原子质量差 Q=$10.6639$ MeV. 这是基态到基态的总可用能量, 末态各动能不唯一, 满足 $T_{\rm Li}+T_e+E_\nu=Q$. 实际可能到 $^8$Li 激发态, 对该支路可用能量为 Q 减激发能, 还可能后续 beta 或粒子发射. 不把激光冷却后的静止母态误解为电子只能有一条固定能线.
\end{solution}
\begin{exercise}\label{ex:2202-9-3}
原题 3. $^{64}$Cu 衰变中观测到中微子和反中微子, (a) 两种 beta 过程、产物是什么? (b) 算各 Q.
\end{exercise}
\begin{solution}
$^{64}$Cu$\to{}^{64}$Zn+$e^-+\bar\nu_e$, Q=$0.57960$ MeV;
$^{64}$Cu$\to{}^{64}$Ni+$e^++\nu_e$, Q=$0.65262$ MeV. 第二式原子质量差须减 $2m_e$. 电子俘获到 Ni 也产 $\nu_e$, Q$_{\rm EC}\simeq1.67462$ MeV, 但它不是题目指定的第二条发射 beta 支路.
\end{solution}
\begin{exercise}\label{ex:2202-9-4}
原题 4. 小 Z 近似忽略 Fermi 函数, beta 谱能量因子为 $\sqrt{T_e^2+2m_eT_e}(T_e+m_e)(Q-T_e)^2$. 对 Q=0.5 MeV, (a) 算并画电子谱; (b) 算并画中微子谱.
\end{exercise}
\begin{solution}
令 $f(T)=\sqrt{T(T+2m_e)}(T+m_e)(Q-T)^2$, $m_e=0.511$ MeV. 电子归一谱为 $f(T)/\int_0^Qf(t)\dd t$, 中微子谱为 $f(Q-E_\nu)$ 除同一积分, 因 $|\dd T/\dd E_\nu|=1$. 两谱镜像, 一般均不是两边对称的抛物线. 下图纵轴使用共同任意常数, 保留真实形状:
\begin{center}\begin{tikzpicture}\begin{axis}[width=9cm,height=5cm,xlabel={动能或 $E_\nu$ (MeV)},ylabel={谱形状 (共同任意单位)},domain=0:0.5,samples=150,legend pos=north east]
\addplot[thick,structurecolor]{sqrt(x*(x+1.022))*(x+0.511)*(0.5-x)^2};\addlegendentry{电子}
\addplot[thick,dashed]{sqrt((0.5-x)*(1.522-x))*(1.011-x)*x^2};\addlegendentry{中微子}
\end{axis}\end{tikzpicture}\end{center}
电子在低动能端因 $p_e$ 趋零, 在高端因中微子相空间趋零. 中微子在零能端因 $E_\nu^2$ 趋零; 在最大端对应电子静止, 同样趋零.
\end{solution}
\begin{exercise}\label{ex:2202-9-5}
原题 5. 允许 beta 过程中电子/正电子和中微子的轨道 $l=0$, 两者自旋 $1/2$. (a) 母核 $I_i$ 时子核可能自旋? (b) 核宇称如何?
\end{exercise}
\begin{solution}
轻子自旋可合成0或1, 故允许 $I_f=I_i$ 或 $I_i\pm1$ (须非负). Fermi 为 $\Delta I=0$, Gamow--Teller 为 $0,\pm1$ 但排除 $0\to0$. 轨道 l=0 不改变核宇称, $\pi_f=\pi_i$. 这是核跃迁宇称选择, 不意味着整个弱作用恢复宇称对称; 轻子手征结构仍破坏 P.
\end{solution}
\begin{exercise}\label{ex:2202-9-6}
原题 6. 中子与静止质量 Am 的核弹性碰撞, 质心散射角 $\theta$, 已知 $E_1/E_0=(A^2+1+2A\cos\theta)/(A+1)^2$. (a) 证明实验室角 $\cos\phi=(1+A\cos\theta)/\sqrt{A^2+1+2A\cos\theta}$. (b) 何时可后散射, 哪种核素不会? (c) 质心各向同性 $P(\cos\theta)=1/2$, 求 $\langle\cos\phi\rangle$.
\end{exercise}
\begin{solution}
(a) 质心平移速度沿原方向 $v_0/(A+1)$, 相对末速为 $Av_0/(A+1)$, x 投影与矢量长度之比正是所求. (b) 后散射需 $1+A\cos\theta<0$, 因而 A>1 且 $\cos\theta<-1/A$. 氢 A=1 的末中子不会严格向后, 180度质心对应中子静止、实验室角不定义. (c) 积分
\begin{equation}
 \langle\cos\phi\rangle=\frac12\int_{-1}^1\frac{1+Au}{\sqrt{A^2+1+2Au}}\dd u.
\end{equation}
置 $w=\sqrt{A^2+1+2Au}$, 则 $(1+Au)=(w^2+1-A^2)/2$, $\dd u=w\dd w/A$, 上下限 $|A-1|,A+1$. 对 A$\ge1$ 求多项式积分得 $2/(3A)$. 对 A=1 为2/3, 重核极限趋0. 若推广到 $0<A<1$, 结果为 $1-A^2/3$, 与 A=1 接合; 核慢化通常只需 A$\ge1$.
\end{solution}
